% GENERATED, edit the French unit or the translation table
\mxnewpage\section*{EV2-T2 -- Mock exam 1}\label{EV2-T2}
\addcontentsline{toc}{section}{EV2-T2 -- Mock exam 1}
\textbf{Practice paper}. Chapters 1 to 7.\par Duration: 2 h. No calculator. Marked out of 20. 14--16 January 2027.
\subsection*{Part 1 -- Fluency (6 points)}
Twelve questions, \num{0,5} point each. Only one correct answer; no justification required. An incorrect or missing answer scores 0, with no loss of points.
\begin{enumerate}
\item Calculate $3+2\times3$.\par A. $12$\quad B. $8$\quad C. $15$\quad D. $9$
\item Calculate $\frac12+\frac13$.\par A. $\frac{2}{5}$\quad B. $\frac{1}{6}$\quad C. $\frac{5}{3}$\quad D. $\frac{5}{6}$
\item Calculate $\dfrac{2^3}{2^{1}}$.\par A. $4$\quad B. $2$\quad C. $8$\quad D. $1$
\item A value of 50 increases by \pc{20}. Final value?\par A. $40$\quad B. $70$\quad C. $60$\quad D. $10$
\item A value of 80 is obtained after a \pc{20} decrease. Initial value?\par A. $64$\quad B. $96$\quad C. $100$\quad D. $60$
\item A \pc{25} increase then a \pc{20} decrease: overall coefficient?\par A. $\frac{9}{10}$\quad B. $\frac{21}{20}$\quad C. $\frac{4}{5}$\quad D. $1$
\item The solution of $2x+3=0$ is\par A. $6$\quad B. $-6$\quad C. $\frac{3}{2}$\quad D. $\frac{-3}{2}$
\item The slope between $\pt{0;3}$ and $\pt{2;9}$ is\par A. $-3$\quad B. $2$\quad C. $6$\quad D. $3$
\item The mean of $3, 5, 7$ is\par A. $3$\quad B. $5$\quad C. $6$\quad D. $15$
\item A sector of $90^\circ$ represents what fraction of the total?\par A. $\frac{1}{4}$\quad B. $\frac{1}{3}$\quad C. $\frac{3}{4}$\quad D. $\frac{1}{2}$
\item A probability is $\dfrac{2}{5}$. Its complement is\par A. $\frac{2}{5}$\quad B. $\frac{5}{2}$\quad C. $\frac{1}{5}$\quad D. $\frac{3}{5}$
\item For $f(x)=x^2+3$, $f(-2)$ is\par A. $12$\quad B. $7$\quad C. $-1$\quad D. $5$
\end{enumerate}
\mxnewpage\subsection*{Part 2 -- Independent exercises (14 points)}
Every step must be justified. Exact results are expected unless stated otherwise.
\mxneed{12}
\subsection*{Exercise 1 -- Two draws (5 points)}
An urn contains four red balls and two blue balls, indistinguishable by touch. Two balls are drawn in succession at random, without replacing the first. At each draw, the balls still present are equally likely. The third question proposes another protocol, with replacement.
\begin{enumerate}
\item Calculate the probability of drawing two red balls. \hfill(1 pt)
\item Calculate the probability of drawing one ball of each colour, taking into account the two possible orders. \hfill(2 pts)
\item The first ball is now replaced before the second draw. Calculate the probability of drawing one ball of each colour and compare it with that of the protocol without replacement. \hfill(2 pts)
\end{enumerate}
\mxneed{12}
\subsection*{Exercise 2 -- A box (5 points)}
Squares of side $x$ cm are cut from the four corners of a square sheet of side 10 cm. Folding up the edges forms a box without a lid. Its volume, in cubic centimetres, is $V(x)=x(10-2x)^2$. We study this model on $\intff{0}{5}$, endpoints included even though they are degenerate.
\begin{enumerate}
\item Calculate $V\prime$ in factorised form. \hfill(1 pt)
\item Study the variation. \hfill(2 pts)
\item Determine the exact maximum volume and justify that the endpoints do not give the maximum. \hfill(2 pts)
\end{enumerate}
\mxneed{12}
\subsection*{Exercise 3 -- Recognition and intersection (4 points)}
In an orthonormal frame, consider the set $\mathcal C$ of points $M\pt{x;y}$ satisfying $x^2+y^2+4x-2y-4=0$. After identifying it, we will study its intersections with a line and a tangent.
\begin{enumerate}
\item Show that $\mathcal C$ is a circle and give its centre and radius. \hfill(1 pt)
\item Determine all the points common to $\mathcal C$ and the line $x=-2$. \hfill(1 pt)
\item Write the equation of the tangent to the circle at the intersection point with the greater y-coordinate. Justify its direction. \hfill(2 pts)
\end{enumerate}
